\section{The degree, level-prime and ramified estimates for the period error}

% BEGIN H88E1
\paragraph{H88E1. Weight and scale inequalities on the original member.}
Fix a family in the design class of X87S4, $0<\epsilon<1/52$,
$x\ge2$, $x>\mathcal F.x_1$, and an index of that family. Use
$q_x=x^\epsilon$, $b=\exp(q_x)$, $S=L.\mathrm{Sc}$,
$\eta=L.\mathrm{etaSz}$, $N_i=k+i(t+1)$, $N=N_m$,
$f=2n$, $h=\#TT$, and $\delta=b'/(xn)$. The word and lift
$1\le j\le m$ are arbitrary.
The denominator of the original $\ell_*$
(the quantity $\mathrm{szlStar}$) is at least
$kn^2(S+t+q_x)\ge km$ and its numerator is $mb$.
The size inequality $b'+S+\eta+b/\ell_*<\ell_*$
therefore gives $kb'\le b$. The parameter bounds give
$kS,m(t+1)\eta\le\delta\le b'/2$.
The cubic lower bound for $\eta/m$ gives $N^2\le\eta$, and
$\eta\le b/(t^3fk^3)$ gives $f\eta\le b$.
Thus $kN\le N^2\le b$ and $5N\le b$, since $k\ge5$.

Write $w=w_j(\tau)=e_{j-1}(\tau)$. The bottom-weight bound gives
$b_s\le(1+\alpha)b'$, where $\alpha=\mathrm{szNst}\,x^\epsilon
L.\mathrm{LaG}/k$. The $\rho$ bound gives $4\alpha\le x^{-2}\le1$
after dropping a nonnegative summand. On top embeddings the window is
$b_s+b'/x\le w\le b_s+b'/x+\eta$, with $b_s\ge b'$.
On thin embeddings it is $\delta\le w\le\delta+\eta$, and on
thick embeddings it is $b\le w\le b+\eta+1$. In particular
$0\le w\le3b$.
The embedding $\tau_1(s)$ is never thick and is top exactly when
$s\in TT$. Thus $w^-=w_j(\tau_1(s))\le9b'/4$.
The other embedding $\tau_2(s)$ is never top, so
$w^+=w_j(\tau_2(s))\le b+\eta+1$. These descriptions follow
from the opposite first coordinates of $\tau_1,\tau_2$ and
injectivity of the second coordinate on each half.
% END H88E1

% BEGIN H88E2
\paragraph{H88E2. Degree count with the original constant 40.}
Set $i=j-1$. At each embedding write $U_r$ for the member's
word-dependent Harish--Chandra head entries and $e_q$ for its
common block parameters. In the sum $\nu_{u,j,\tau}$ of H87P2
the first $k$ entries are
\[
 z_r=U_r+w-(N_i+t/2-r)\qquad(0\le r<k).
\]
The subtracted quantity is in $[0,N]$. For each $q<i$ there are
$t+1$ further equal entries $w-e_q-(i-q)(t+1)$.
They are positive by the spacing of the block parameters and
at most $2\eta$ because $|e_q-b_0(\tau)|,|w-b_0(\tau)|\le\eta$.
The middle $k-2$ head entries have $0\le U_r\le S$.

Put $a_\tau=a_s$ at $\tau=\tau_1(s)$ with $s\in TT$, and zero
elsewhere. At a top embedding,
$U_0=a_s+(k-1)/2$, $U_{k-1}=-b_s-(k-1)/2$. Hence
$z_0+z_{k-1}-w-a_s$ has absolute value at most
$2N+b'/x+\eta$, by the top window in H88E1. At other embeddings
$0\le U_0\le3S$ by interlacing with
$0\le v_0\le S$ and $U_0-v_0\le2S$.
Also $U_{k-1}=v_{k-2}-1/2$ and $0\le v_{k-2}\le S$.
Here $|z_0+z_{k-1}-2w|\le4S+2N+1$. With
$w_{\rm off}=w$ off the tops and zero on them, summing over rows gives
\[
 \begin{split}
 |\nu_{u,j,\tau}-(k-1)w-a_\tau-w_{\rm off}|
 \le {}&k(S+N)+2m(t+1)\eta+4S+2N+1+\eta\\
       &+{\bf1}_{\tau\ {\rm top}}b'/x .
 \end{split}
\]

The Fock degree at this same embedding is the sum of the absolute
values of its $\kappa$ entries. The positive-signature coordinates
have the form
\[
 s_r+w-(n_V-1-2r)/2-(k+j-1-(m-j))/2,
 \qquad 0\le r<k-1+j ,
\]
where $s_r=v_r$ for $r<k-1$ and $s_r=-b_{r-k+1}(\tau)$ thereafter.
The other $m-j$ coordinates have absolute values
$|b_{j+c}(\tau)-w-(t+1+2c)/2|$ for $0\le c<m-j$.
The combined rank shifts have absolute value at most $2N$,
as $n_V=k+m-1+t$, $k+m\le N$ and $t+1\le N$.
For the first $k-1$ entries use
$\big||z|-w\big|\le |z-w|$ and $0\le v_r\le S$.
The other $m$ entries use $|b_q(\tau)-w|\le\eta$.
Writing $d_{j,\tau}$ for this contribution gives
\[
 |d_{j,\tau}-(k-1)w|\le k(S+2N)+2m(\eta+2N).
\]
These are the stored coordinates by the member's $\mathrm{KapDef}$.

The sum of the two error bounds, excluding the top contribution
$b'/x$, is at most $20(kS+N^2+mt\eta)$. Indeed $kN,mN,N,1\le N^2$,
$S\le kS$ and $m\eta\le mt\eta$.
There are $f$ embeddings and each of the $h$ top labels occurs once.
Consequently, with $A=\sum_{s\in TT}a_s$,
\[
 \left|\sum_\tau \nu_{u,j,\tau}-d_j-A-
                         \sum_{\tau\ {\rm off}}w_j(\tau)\right|
 \le20\{f(kS+N^2+mt\eta)+h b'/x\}.
\]
Let $Y=f(kS+mt\eta+kN)+hb'/x$. Since $N^2\le mt\eta$,
the braces are at most $2Y$. Since $w\ge0$, this proves
the source inequality with right side
$\mathrm{errG}(L,40,j)=2\sum_\tau w_j(\tau)+40Y$.
Also $h\le f$ and each of $kS,mt\eta,kN,b'/x$ is at most $b$,
so $Y\le4fb$ and $\mathrm{errG}(L,40,j)\le166fb$.
% END H88E2

% BEGIN H88E3
\paragraph{H88E3. Both distinguished-prime bounds and the stabilizer term.}
On its eventual range $\mathrm{DGamSize}$ supplies a single $C_D$ with
$\log D_{\Gamma,u}\le C_D fN\log(fN)$. The elementary inequality
$\log(fN)\le fN$ and H88E1 give
$fN\log(fN)\le f(fN^2)\le f(f\eta)\le fb$.
Thus $C_\Gamma=2|C_D|$ bounds $2\log D_{\Gamma,u}/B_0$,
where $B_0=fmbq_x$. At either distinguished prime $g$, the parameter
choice has $g\le2Y_g$ with $\log Y_g\le C_{YL}q_x$, hence
$\log g\le(C_{YL}+1)q_x$. Also
$v_g(D_{\Gamma,u})\log g\le\log D_{\Gamma,u}$.
For the remaining error put $r=c.g\,f(N^2+m^2\eta)$ as in H87P2.
Since $m\le f$, $N^2\le\eta$, and $f\eta\le b$, the
quantity $N^2+m^2\eta$ is at most $2mb$. H88E2 therefore gives
\[
 \big(\mathrm{errG}(L,40,j)+2v_g(D_{\Gamma,u})+2r\big)\log g
 \le\big((166+4|c.g|)(C_{YL}+1)+2|C_D|\big)B_0 .
\]
The family fixes $C_{YL}$ and $c.g$ before the choice of member.
This gives the same $C_g$ at both distinguished primes.
% END H88E3

% BEGIN H88E4
\paragraph{H88E4. The error at the word's level prime.}
In the level-prime formula the exceptional first entry at a top
embedding is removed. The remaining head parameters $U_r$ have
absolute value at most $3b$. For interior entries this is the cloud
bound $0\le U_r\le S$. For first entries off the tops use H88E2.
For the last entry on the tops use
$U_{k-1}=-b_s-(k-1)/2$, $b_s\le5b'/4$, $k\le b$.
Off the tops this entry equals $v_{k-2}-1/2$.
In the stored dominant weight the head has the further shift
$-(k-1-2r)/2$. The block with index $1\le a\le i'$ has
weight $-e_{a-1}+k/2+(2a-1)(t+1)/2$ at level $i'\le m$.
For both kinds of entries the definition of $W_\tau$ also subtracts
$\lfloor i'(t+1)/2\rfloor$. In the head these shifts together have
absolute value at most $2N$. In the block the difference of the
rank shifts has absolute value at most $2N$. Since $e_{a-1}\in[0,3b]$
and $5N\le b$, every retained absolute value is at most $5b$.

Thus the nonnegative value $E_T$ at a seen place,
\[
            2\sum_\tau W_\tau+fN+
                \#\{\tau\ {\rm top}\}(t+1)(m+1)/4 ,
\]
is at most $12fb$. Here $(t+1)(m+1)\le2N\le b$.
At an unseen place the value is zero. For the precise mass of the
seen places, let $y$ be the image of the chosen integral generator
$\pi_T\in\mathbb Q(\sqrt{-d})$ with $y\bar y=T=T_u$.
All retained embeddings restrict identically to this quadratic
field, and the stored criterion for being seen says $v_w(y)>0$.
The trace $a=y+\bar y$ is an integer with $a^2\le4T$.
The generator is $1$ modulo $4$ in the integers of the host
field, so $a\equiv2\pmod4$ and $T\equiv1\pmod4$
by integrality and contraction to $\mathbb Z$. Consequently
$T\ge5$, $a\ne0$, and $|a|<T$. A prime over $T$ cannot divide
both $y$ and $\bar y$, since it would then divide $a$.
Their two nonnegative valuations sum to $e_w$. Hence at a seen
place $v_w(y)=e_w$, and elsewhere it is zero. The norm identity
in H87P2 gives $2\sum_{w\mid T,\ {\rm seen}}\ell_we_w=\log T$.
Any subset of the fibre contributes at most $12fb\log T$.
The eventual $\mathrm{TuSize}$ inequality $\log T\le C_{Tu}mq_x$
therefore gives $C_T=\max(12C_{Tu},0)$.
% END H88E4

% BEGIN H88E5
\paragraph{H88E5. The ramified correction and the original eighth entry.}
For $s\in HH$ use the two embeddings $\tau_1(s)$ and $\tau_2(s)$
of H88E1. The term $c_s$ entering the ramified supremum is the
sum of all row entries at $\tau_1(s)$ when $s\notin TT$. At $s\in TT$
remove its first row and add one extra copy of its second row
and the first row at $\tau_2(s)$. This is exactly
$\mathrm{ctD}$, with the arbitrarily large first top weight cancelled.
The remaining head at $\tau_1(s)$ has $k$ terms in either case,
counting multiplicity. Each is at most $w^-+3S$, since the cloud
bounds apply to the interior, the last top parameter is negative,
and the first off-top parameter is at most $3S$. The tail contributes
at most $2m(t+1)\eta$ by H88E2. The extra first row on the other
half is at most $w^++3S\le2b$. Indeed $\eta,3S\le b'/2$,
$kb'\le b$, $k\ge5$ and $5N\le b$ imply $1+\eta+3S\le b$.
All these row terms are nonnegative. Thus
\[
 0\le c_s\le 2b\,{\bf1}_{s\in TT}+
      k(9b'/4)+3kS+2m(t+1)\eta
 \le 2b\,{\bf1}_{s\in TT}+4kb'.
\]
There are $n=f/2$ labels. Taking the supremum over the nonempty
word set bounds its half by $hb+fkb'$ from above and zero from below.
The other term in $E_{d,j}$ is
$-(k-1)\sum_s\min(w^-,w^+)/2$. It lies in $[-fkb',0]$
by $w^-\le9b'/4$ and $n=f/2$.
Consequently $|E_{d,j}|\le hb+fkb'\le2fb$. The
member's discriminant-size inequality gives
$\log d\le x^{\epsilon/5+\epsilon/100}\le q_x$. It follows that
$2\log d\,|E_{d,j}|\le4B_0$, so $C_d=4$.

Combine H88E2--H88E5 with the two inert-prime estimates I88P1--I88P3
and the fibre calculation H87P3. For every family in the design
class, and every $0<\epsilon<1/52$, this proves the entire eighth
entry on the actual members on one eventual range. One may use
\[
 C=2C_g+\max(12C_{Tu},0)+2|C_D|+4+
                         18(1+1/\epsilon)+848 .
\]
Take the maximum of $2,\mathcal F.x_1+1$ and the two thresholds
for $\mathrm{TuSize}$ and $\mathrm{DGamSize}$, together with the
threshold of I88P3. Every place, word and lift in the original
error sum is covered. The other entries of the period-height
contract retain their own hypotheses.
% END H88E5
